\documentclass[12pt]{article}
\usepackage[T1]{fontenc}
\usepackage[utf8]{inputenc}
\usepackage{lmodern}
\usepackage{textcomp}
\usepackage{enumitem}
\usepackage{geometry}
\usepackage{graphicx}
\usepackage{amsmath, amssymb}
\geometry{margin=1in}
\begin{document}
\begin{center}
Physics - Undergraduate Level Assessment 8 \
Total Marks: 40 \
Answer all questions.
\end{center}

\section*{Multiple Choice (10 marks)}
\begin{enumerate}[label=\arabic*.]
\item The emission spectrum of the doubly ionized lithium atom Li++ (Z = 3, A = 7) is identical to that of a hydrogen atom in which all the wavelengths are
\begin{enumerate}[label=(\alph*)]
    \item decreased by a factor of 9
    \item decreased by a factor of 49
    \item decreased by a factor of 81
    \item increased by a factor of 9
\end{enumerate}
\item The rest mass of a particle with total energy 5.0 GeV and momentum 4.9 GeV/c is approximately
\begin{enumerate}[label=(\alph*)]
    \item 0.1 GeV/$c^2$
    \item 0.2 GeV/$c^2$
    \item 0.5 GeV/$c^2$
    \item 1.0 GeV/$c^2$
\end{enumerate}
\item Which of the following lasers utilizes transitions that involve the energy levels of free atoms?
\begin{enumerate}[label=(\alph*)]
    \item Diode laser
    \item Dye laser
    \item Free-electron laser
    \item Gas laser
\end{enumerate}
\item For which of the following thermodynamic processes is the increase in the internal energy of an ideal gas equal to the heat added to the gas?
\begin{enumerate}[label=(\alph*)]
    \item Constant temperature
    \item Constant volume
    \item Constant pressure
    \item Adiabatic
\end{enumerate}
\item Which of the following gives the total spin quantum number of the electrons in the ground state of neutral nitrogen (Z = 7)?
\begin{enumerate}[label=(\alph*)]
    \item 1/2
    \item 1
    \item 3/2
    \item 5/2
\end{enumerate}
\end{enumerate}

\section*{True / False (10 marks)}
\textit{Answer True or False}
\begin{enumerate}[label=\arabic*.]
\setcounter{enumi}{5}
\item True or False: If the total energy of a particle is twice its rest energy, its relativistic momentum is equal to (\ensuremath{\sqrt} 3)mc.
\item True or False: Electrical resistivity measurements can directly indicate the type of charge carriers in a semiconductor.
\item True or False: In a constant pressure process, the increase in the internal energy of an ideal gas is equal to the heat added to the gas.
\item True or False: A meter stick moving at 0.8 c takes 1.6 ns to pass an observer in their reference frame.
\item True or False: During a reversible thermodynamic process, the internal energy of the system always increases.
\end{enumerate}

\section*{Long Form Response (20 marks)}
\textit{Answer each question with a clear and complete explanation.}
\begin{enumerate}[label=\arabic*.]
\setcounter{enumi}{10}
\item Discuss the work required to accelerate a proton from rest to a speed of 0.6 c. Explain the underlying physics principles and derive the expression for the work done in terms of the proton's mass.
\item Discuss the measurements an astronomer can make when observing a moon orbiting a planet, and explain why the mass of the moon cannot be determined from these measurements.
\end{enumerate}

\vfill
\noindent\textit{End of Paper}
\end{document}